The characterization in Section~\ref{Results} maps the operating envelope of the composed system honestly. We read each apparent limitation as a quantified property that an operator can plan around, rather than a defect to be hidden.

\paragraph{Bias absorption is a property of the interpretable model, not a bug.}
The monotone dependence of bias recall on conditional standard deviation (Table~\ref{tab:biascoupling}) follows directly from what makes the base detector interpretable. Because the conditional mean of a strongly coupled signal tracks its peers, a constant offset injected on that signal is partly explained away by the conditional model, and the residual---hence the typing---does not see it. The same mechanism that yields tight, peer-informed operating limits also absorbs offsets on those signals. The practical reading is that bias is reliably typed where it is observable (weakly coupled signals show recall up to $0.941$) and is absorbed where the model is confident enough to overrule it; the aggregate $0.518$ is the average of these two regimes, not a uniform failure. An operator who needs bias coverage on a strongly coupled signal should monitor that signal's raw calibration directly, a check the conditional model is, by construction, not designed to provide.

\paragraph{Two-signal cancellation limits the real-deployment false-positive floor.}
On the Intel Lab stream the residual healthy false positives are $\approx 93\%$ drift, traced to diurnal ramps that a two-signal conditional model cannot fully cancel. This is the same correlation-dependence quantified in Table~\ref{tab:regime}: with few correlated signals, slow shared movement leaks into the residual. The limitation is therefore structural and predictable---it improves with the number of correlated signals available---rather than a tuning failure. We report the $7.4\%$ floor at $90\%$ detection as the achievable operating point on a two-signal deployment, with the explicit expectation that richer instrumentation would lower it.

\paragraph{The dead-band is the priced cost of adaptation robustness.}
Because robustness to faults arriving during adaptation comes from graded change-point suppression and not from residual cancellation (Section~\ref{sec:regime}), it carries a cost: the masking dead-band of roughly $[3\sigma, 15\sigma]$ for faults co-occurring with a change point (Table~\ref{tab:deadband}). We present this as an explicit trade-off governed by a single parameter, \texttt{suppress\_scale}. A deployment that fears spurious labels from coordinated shifts should keep the default suppression; one that fears missing a fault during a regime change should lower it, accepting the false-label rate of Table~\ref{tab:regime}. The point is that the trade-off is measured and tunable, not silent.

\paragraph{Threshold and freeze-constant calibration are deployment decisions.}
The operating-point sweep (Table~\ref{tab:operating}) and the freeze-on-quiescent ablation (Table~\ref{tab:quiescent}) both show that the system's behavior is set by a small number of interpretable constants. The mean threshold selects the false-alarm/detection trade-off; the freeze constants determine how aggressively a quiet signal is judged frozen, with the strict setting eliminating false freeze labels on a live-but-quiet signal that the default mislabels $72\%$ of the time. We recommend the strict freeze constants or per-signal calibration for steady-state signals, and we expose the threshold as the operator's primary control. These are calibration decisions with quantified consequences, not hidden defaults.

\paragraph{Data availability constrains the evaluation to triangulation.}
No public stream simultaneously labels all four fault types per signal and per timestep, which is why we combine controlled injection (per-type ground truth, but synthetic) with a real fault deployment (genuine faults, but a single physical failure mode with a held-out label). Neither alone would be conclusive: the injections could overstate performance on idealized faults, and the real deployment exercises mainly the freezing signature. Read together, the injections establish the per-type envelope and the refuted claims, while the real deployment confirms that the envelope holds on genuine faults at a usable operating point. We regard this triangulation as the appropriate evidentiary standard given the data that exist, and we flag the absence of a fully labelled multi-type stream as the principal external constraint on stronger claims.
